% Ch1_2, slide 10: the parallelepiped of A, B and C; its base has area |B x C|
\begin{tikzpicture}[scale=1]
  \coordinate (O) at (0,0); \coordinate (Z) at (0,3);
  \coordinate (B) at (3.6,-0.7); \coordinate (C) at (1.5,0.6); \coordinate (A) at (0.9,2.1);
  \fill[ELfillB] (O)--(B)--($(B)+(C)$)--(C)--cycle;
  \draw[ELdash] (C)--($(B)+(C)$)--(B);
  \draw[ELdash] (A)--($(A)+(C)$)--($(A)+(B)+(C)$)--($(A)+(B)$)--cycle;
  \draw[ELdash] (C)--($(A)+(C)$) ($(B)+(C)$)--($(A)+(B)+(C)$) (B)--($(A)+(B)$);
  \draw[ELvec] (O)--(B) node[ELlab,anchor=north]{$\vect{B}$};
  \draw[ELvec] (O)--(C) node[ELlab,anchor=south east,xshift=2pt]{$\vect{C}$};
  \draw[ELvec] (O)--(A) node[ELlab,anchor=east]{$\vect{A}$};
  \draw[ELvec2] (O)--(0,3.2) node[ELlab,anchor=south]{$\vect{B}\times\vect{C}$};
  \pic[draw=ELink!70,line width=0.5pt,angle radius=0.55cm,"$\theta_1$",angle eccentricity=1.55,font=\footnotesize] {angle=B--O--C};
  \pic[draw=ELink!70,line width=0.5pt,angle radius=0.75cm,"$\theta_2$",angle eccentricity=1.4,font=\footnotesize] {angle=A--O--Z};
  \node[ELlab,anchor=north] at (1.6,-0.75) {\ELL{Area}{مساحت} $=\abs{\vect{B}\times\vect{C}}$};
\end{tikzpicture}
